\documentclass[10pt,a4paper]{amsart}
\usepackage[margin=19mm]{geometry}
\usepackage{amsmath,amssymb,mathtools}
\usepackage[scr=rsfs,cal=euler]{mathalfa}
\usepackage{xfrac}
\usepackage[unicode,hidelinks,bookmarks=false]{hyperref}
\newcommand{\ie}{i.e.\ }
\newcommand{\cf}{cf.\ }
\newcommand{\resp}{respectively\ }
\newcommand{\Subsection}{\subsection}
\providecommand{\nobreakdash}{\nobreak\hbox{-}}
\setlength{\emergencystretch}{2em}
\multlinegap0pt
\usepackage{amssymb}
\usepackage[shortlabels]{enumitem}
\usepackage{mathtools}
\usepackage{xcolor}
\usepackage{tikz}
\newtheorem{theorem}{Theorem}
\newtheorem{lemma}{Lemma}
\newcommand{\F}{\mathbb F}
\newcommand{\Z}{\mathbb Z}
\newcommand{\balpha}{\boldsymbol{\alpha}}
\newcommand{\cL}{\mathcal L}
\title{Nine-distance theorem and growth of best-approximation denominators}
\author{Nikita Shulga}
\date{Source: arXiv:2608.01443v1 (CC BY 4.0)}
\begin{document}
\maketitle

\noindent\emph{Source and licence.} Nine-distance theorem and growth of best-approximation denominators. Version arXiv:2608.01443v1 (CC BY 4.0), distributed by arXiv under the Creative Commons Attribution 4.0 International licence, http://creativecommons.org/licenses/by/4.0/, as recorded in the arXiv licence metadata for this exact version and shown on the abstract page https://arxiv.org/abs/2608.01443v1. Full licence notice: \texttt{licenses/arxiv-2608.01443-CC-BY-4.0.txt}.\par\smallskip

\noindent\textit{Editorial scope.} Counted unit: the article's theorem thm:growth (source0.tex lines 530--550). The article's complete original proof of it is delivered with it (source0.tex lines 680--871, one \texttt{proof} environment of 192 lines, which the article places away from the statement). No mathematics is written, repaired or re-derived: the statement and the proof are the article's own lines, in the article's own notation, and no retained line is substituted or re-flowed.  Retained uncounted as the source-local context the proof needs: lines 615--620; lines 633--641; lines 665--672.  Nothing was omitted from any retained span, so the package carries no omission record.  No retained line contains a citation, so the wrapper carries no bibliography and no bibliography entry.  No retained line contains a figure environment, an image-inclusion command or drawing code, so no image is delivered and the package declares no asset dependency.  Editorial formatting only, in addition: an AMS \texttt{amsart} preamble that restates the article's own declarations and macros used by the retained lines, namely the environments \texttt{theorem}, \texttt{lemma} on separate counters, together with the standard packages those lines need.  The retained environments print in this document as Theorem 1, Lemma 1 in order of appearance, whereas the article prints the same items with the numbering of its own full document.  The complete native source of the article as distributed in the arXiv e-print for this exact version is bundled unchanged as \texttt{sources/206629/source0.tex}.
\begin{lemma}\label{lem:record}
For every defined best denominator $q_n$,
$$
 1\le q<q_n \quad\Longrightarrow\quad \rho(q)>r_n.
$$
\end{lemma}

\begin{lemma}\label{lem:sum-free}
Let $m\ge1$ and let $A\subset \F_2^m\setminus\{0\}$. If no three distinct
elements of $A$ have sum zero, then
$$
 |A|\le 2^{m-1}.
$$
If equality holds, then $A$ is a nonzero coset of a codimension-one subspace
of $\F_2^m$.
\end{lemma}

\begin{lemma}\label{lem:four-sign}
For all $a,b,c$ in a real inner-product space,
\begin{align*}
 &\|a+b-c\|^2+\|a-b+c\|^2
   +\|-a+b+c\|^2+\|a+b+c\|^2 \\
 &\hspace{42mm}=4\bigl(\|a\|^2+\|b\|^2+\|c\|^2\bigr).
\end{align*}
\end{lemma}

\begin{theorem}\label{thm:growth}
Put $M=2^d$. Whenever $q_{n+M}$ is defined, one of the following alternatives
holds:
\begin{enumerate}
\item
$$
 q_{n+M}\ge2q_{n+1};
$$
\item the set $\{1,\ldots,M\}$ can be partitioned into disjoint pairs $\{j,k\}$,
$j<k$, such that
\begin{equation}\label{eq:pairing-recurrence}
 q_{n+k}=q_n+q_{n+j}.
\end{equation}
\end{enumerate}
Consequently, we always have
\begin{equation}\label{eq:refined-growth}
 q_{n+M}\ge
 \min\{2q_{n+1},\,q_n+q_{n+M/2}\}
 \ge q_n+q_{n+1}>2q_n.
\end{equation}
\end{theorem}

\begin{proof}[Proof of Theorem~\ref{thm:growth}]
Put $M=2^d$ and write
$$
 Q_i=q_{n+i}\qquad(0\le i\le M).
$$
If $Q_M\ge2Q_1$, the first alternative holds, so suppose that
\begin{equation}\label{eq:exceptional-block}
 Q_M<2Q_1.
\end{equation}
For each $i$, choose $P_i\in\cL$ such that
$$
 x_i=Q_i\balpha-P_i,
 \qquad \|x_i\|=\rho(Q_i)=:r_i,
$$
and put $R=r_0$. Thus
\begin{equation}\label{eq:errors}
 R=r_0>r_1>\cdots>r_M.
\end{equation}
Since $Q_1$ is the first denominator after $Q_0$ at which the approximation
error is strictly smaller than $R$, we have
\begin{equation}\label{eq:below-Q1}
 0<q<Q_1\quad\Longrightarrow\quad \rho(q)\ge R.
\end{equation}

Choose a basis of $\cL$, and let $p_i\in\Z^d$ be the coordinate vector of
$P_i$. Set
$$
 z_i=(Q_i,p_i)\in\Z^{d+1},
 \qquad b_i=z_i\pmod 2\in\F_2^{d+1}.
$$
We first show that
$$
 \mathcal A=\{b_1,\ldots,b_M\}
$$
is a sum-free subset of $\F_2^{d+1}$ of cardinality $M$.

\smallskip
\noindent\emph{The classes are nonzero.}
If $b_i=0$ for some $1\le i\le M$, then $z_i/2$ is integral and
$$
 0<\frac{Q_i}{2}<Q_1
$$
by \eqref{eq:exceptional-block}. The vector $x_i/2$ is admissible in the
definition of $\rho(Q_i/2)$, and therefore
$$
 \rho(Q_i/2)\le\frac{r_i}{2}<R,
$$
contrary to \eqref{eq:below-Q1}.

\smallskip
\noindent\emph{The classes are distinct.}
Suppose $b_i=b_j$ with $1\le i<j\le M$. Then $(z_j-z_i)/2$ is integral and
$$
 0<\frac{Q_j-Q_i}{2}<Q_1.
$$
Consequently,
$$
 \rho\left(\frac{Q_j-Q_i}{2}\right)
 \le\frac12\|x_j-x_i\|
 \le\frac{r_i+r_j}{2}<R,
$$
again contradicting \eqref{eq:below-Q1}.

\smallskip
\noindent\emph{No three classes sum to zero.}
Suppose, to the contrary, that
\begin{equation}\label{eq:zero-sum-triple}
 b_i+b_j+b_k=0,
 \qquad 1\le i<j<k\le M.
\end{equation}
Define
$$
 \begin{aligned}
 A&=x_i+x_j-x_k,&
 B&=x_i-x_j+x_k,\\
 C&=-x_i+x_j+x_k,&
 D&=x_i+x_j+x_k.
 \end{aligned}
$$
The four analogous signed combinations of $z_i,z_j,z_k$ lie in
$2\Z^{d+1}$. Moreover,
$$
 0<\frac{Q_i+Q_j-Q_k}{2}<Q_1,
 \qquad
 0<\frac{Q_i-Q_j+Q_k}{2}<Q_1.
$$
Indeed, the lower bound in the first inequality follows from
$Q_i\ge Q_1$, $Q_j>Q_1$, and $Q_k<2Q_1$; all the other bounds follow from
the ordering of the $Q_i$ and \eqref{eq:exceptional-block}. Hence
\eqref{eq:below-Q1} gives
\begin{equation}\label{eq:AB-long}
 \|A\|\ge2R,\qquad \|B\|\ge2R.
\end{equation}

Lemma~\ref{lem:four-sign}, applied to $x_i,x_j,x_k$, gives
$$
 \|A\|^2+\|B\|^2+\|C\|^2+\|D\|^2
 =4\bigl(r_i^2+r_j^2+r_k^2\bigr).
$$
Using \eqref{eq:AB-long} and $\|D\|^2\ge0$, we obtain
$$
 \begin{aligned}
 \|C\|^2
 &\le4\bigl(r_i^2+r_j^2+r_k^2\bigr)-8R^2\\
 &=4r_k^2+4\bigl(r_i^2+r_j^2-2R^2\bigr)
 <4r_k^2,
 \end{aligned}
$$
because $i,j\ge1$ and hence $r_i,r_j<R$. Thus
\begin{equation}\label{eq:C-short}
 \|C\|<2r_k.
\end{equation}
On the other hand, $(-z_i+z_j+z_k)/2$ is integral and has denominator
$$
 t=\frac{-Q_i+Q_j+Q_k}{2},
 \qquad 0<t<Q_k.
$$
Therefore
$$
 \rho(t)\le\frac{\|C\|}{2}<r_k=\rho(Q_k),
$$
contrary to Lemma~\ref{lem:record}. This proves that $\mathcal A$ is
sum-free.

Since $|\mathcal A|=M=2^d=\frac12|\F_2^{d+1}|$,
Lemma~\ref{lem:sum-free} gives
\begin{equation}\label{eq:affine-hyperplane}
 \mathcal A=a+H,
\end{equation}
where $H$ is a codimension-one subspace and $a\notin H$.

We next show that $b_0$ is a nonzero element of $H$. If $b_0=0$, then
$Q_0/2<Q_1$ and
$$
 \rho(Q_0/2)\le R/2<R,
$$
contrary to \eqref{eq:below-Q1}. If $b_0=b_i$ for some $1\le i\le M$, then
$$
 0<\frac{Q_i-Q_0}{2}<Q_1
$$
and
$$
 \rho\left(\frac{Q_i-Q_0}{2}\right)
 \le\frac{r_i+R}{2}<R,
$$
again a contradiction. Thus $b_0\notin\mathcal A$ and $b_0\ne0$, so
\eqref{eq:affine-hyperplane} implies $b_0\in H$.

Translation by $b_0$ is therefore a fixed-point-free involution of
$\mathcal A$. It partitions $\{1,\ldots,M\}$ into pairs $\{j,k\}$ such that,
after ordering each pair with $j<k$,
\begin{equation}\label{eq:parity-pair}
 b_k=b_0+b_j.
\end{equation}
We claim that every such pair satisfies $Q_k=Q_0+Q_j$. Suppose otherwise.
The relation \eqref{eq:parity-pair} implies that the signed combinations of
$z_0,z_j,z_k$ used below are even. The two nonzero denominators
$$
 \frac{|Q_0+Q_j-Q_k|}{2},
 \qquad
 \frac{Q_0-Q_j+Q_k}{2}
$$
both lie in $(0,Q_1)$, by the ordering of the $Q_i$ and
\eqref{eq:exceptional-block}. The corresponding vectors are,
up to sign,
$$
 x_0+x_j-x_k,\qquad x_0-x_j+x_k.
$$
Hence \eqref{eq:below-Q1} shows that both of these vectors have norm at
least $2R$. The four-sign identity, now applied
to $x_0,x_j,x_k$, then gives
$$
 \|-x_0+x_j+x_k\|^2
 \le4(R^2+r_j^2+r_k^2)-8R^2
 <4r_k^2.
$$
But $(-z_0+z_j+z_k)/2$ has a positive denominator smaller than $Q_k$, which
again contradicts Lemma~\ref{lem:record}. Hence
$$
 Q_k=Q_0+Q_j,
$$
proving \eqref{eq:pairing-recurrence}.

Finally, let $\{s,M\}$ be the pair containing $M$. There are $M/2$ smaller
members of the pairs, and $s$ is the largest of them, because the map
$Q_j\mapsto Q_0+Q_j$ preserves order. Hence $s\ge M/2$ and
$$
 Q_M=Q_0+Q_s\ge Q_0+Q_{M/2}.
$$
This proves \eqref{eq:refined-growth}; its final strict inequality follows
from $Q_1>Q_0$ and $Q_{M/2}>Q_0$.
\end{proof}

\end{document}
